\chapter{Consecuencias globales de la derivada}
\label{chap:consecuencias-derivada}

La derivada es local, pero el teorema del valor medio permite convertir su
signo en información sobre todo un intervalo. Esa conexión organiza este
capítulo: primero estudiaremos variación y extremos; después formularemos
problemas de optimización; por último, usaremos la segunda derivada para
describir convexidad e inflexión.

\section{Monotonía y teorema del valor medio}
\label{sec:monotonia-mvt}

\subsection{Cuatro nociones de monotonía}
\CanonicalUnit{CM-C04-U001}{}{}

Sea \(I\subseteq\R\) un intervalo y \(f:I\to\R\). Diremos que \(f\) es

\begin{itemize}
  \item \emph{no decreciente} en \(I\) si
        \(x<y\) implica \(f(x)\le f(y)\);
  \item \emph{no creciente} en \(I\) si
        \(x<y\) implica \(f(x)\ge f(y)\);
  \item \emph{estrictamente creciente} en \(I\) si
        \(x<y\) implica \(f(x)<f(y)\);
  \item \emph{estrictamente decreciente} en \(I\) si
        \(x<y\) implica \(f(x)>f(y)\).
\end{itemize}

Estas son propiedades en un intervalo, no propiedades aisladas de un punto.
Cuando interese comparar el comportamiento a ambos lados de \(c\), se
formularán las desigualdades locales exactas en vez de usar la expresión no
estándar ``creciente en \(c\)''.

La distinción entre desigualdades débiles y estrictas es esencial. Una función
constante es no decreciente y no creciente, pero no es estrictamente monótona
en ningún intervalo no degenerado.

\subsection{Dependencias ya establecidas}
\CanonicalUnit{CM-C04-U002}{}{}

El teorema de Rolle y el teorema del valor medio ya fueron enunciados y
demostrados en la Sección~\ref{sec:rolle-valor-medio}. No se repiten aquí. En
lo que sigue se utilizará, en particular, que para una función continua en
\([a,b]\) y derivable en \((a,b)\) existe \(c\in(a,b)\) tal que

\[
  f(b)-f(a)=f'(c)(b-a).
\]

Esta identidad es el puente entre el signo local de \(f'\) y la comparación
global de valores de \(f\).

\subsection{Criterios de monotonía mediante la derivada}
\CanonicalUnit{CM-C04-U003}{}{}

\begin{teorema}[Signo de la derivada y monotonía]
Sea \(I\) un intervalo, \(f:I\to\R\) continua en \(I\) y derivable en
\(\operatorname{int}I\). Entonces:

\begin{enumerate}
  \item si \(f'(x)\ge0\) para todo \(x\in\operatorname{int}I\), \(f\) es no
        decreciente en \(I\);
  \item si \(f'(x)\le0\) para todo \(x\in\operatorname{int}I\), \(f\) es no
        creciente en \(I\);
  \item si \(f'(x)>0\) para todo \(x\in\operatorname{int}I\), \(f\) es
        estrictamente creciente en \(I\);
  \item si \(f'(x)<0\) para todo \(x\in\operatorname{int}I\), \(f\) es
        estrictamente decreciente en \(I\).
\end{enumerate}
\end{teorema}

\begin{proof}
Tomemos \(x<y\) en \(I\). El teorema del valor medio aplicado a \([x,y]\)
proporciona \(c\in(x,y)\) tal que

\[
  f(y)-f(x)=f'(c)(y-x).
\]

Como \(y-x>0\), el signo de la diferencia \(f(y)-f(x)\) es el de \(f'(c)\).
Las cuatro conclusiones se obtienen usando, respectivamente, una desigualdad
débil o estricta y su versión con signo contrario.
\end{proof}

La conversa estricta es falsa. La función \(f(x)=x^3\) es estrictamente
creciente en \(\R\), pero \(f'(0)=0\). En cambio, si una función no
decreciente es derivable en un punto interior, su derivada en ese punto no
puede ser negativa; la afirmación análoga vale para funciones no crecientes.

\subsection{Derivada nula e igualdad de derivadas}
\CanonicalUnit{CM-C04-U004}{}{}

\begin{corolario}[Derivada nula]
Sea \(I\) un intervalo. Si \(f:I\to\R\) es continua en \(I\), derivable en
\(\operatorname{int}I\) y \(f'(x)=0\) en el interior, entonces \(f\) es
constante en \(I\).
\end{corolario}

\begin{proof}
Para cualesquiera \(x<y\) en \(I\), el teorema del valor medio da
\(f(y)-f(x)=f'(c)(y-x)=0\). Por tanto, todos los valores coinciden.
\end{proof}

Si \(f'=g'\) en el interior de un intervalo y ambas funciones satisfacen las
hipótesis anteriores, entonces \((f-g)'=0\). En consecuencia, existe
\(K\in\R\) tal que

\[
  f(x)=g(x)+K
  \qquad(x\in I).
\]

La conclusión depende de que el dominio sea un intervalo. En un dominio con
varias componentes, la constante podría ser distinta en cada componente.

\subsection{Dos consecuencias representativas}
\CanonicalUnit{CM-C04-U005}{}{}

Si \(f'(x)>0\) en un intervalo, cada ecuación \(f(x)=\lambda\) tiene a lo sumo
una solución allí, porque la estricta monotonía impide repetir valores. La
existencia de una solución exige, además, información de continuidad y de
valores o límites; monotonía no significa existencia.

Como ejemplo de desigualdad, el teorema del valor medio aplicado a la
exponencial permite probar

\[
  e^x\ge 1+x \qquad(x\in\R),
\]

con igualdad solo en \(x=0\). Para \(x>0\), se escribe
\(e^x-1=e^c x\) con \(c\in(0,x)\) y \(e^c>1\). Para \(x<0\), se obtiene la
misma identidad con \(c\in(x,0)\); entonces \(0<e^c<1\) y la multiplicación
por \(x<0\) conserva la desigualdad en la dirección requerida.

\section{Extremos y criterios derivados}
\label{sec:extremos-criterios}

\subsection{Extremos, puntos estacionarios y puntos críticos}
\CanonicalUnit{CM-C04-U006}{}{}
\ProofRole{proof-026}{}

Sea \(f:D\to\R\) y \(c\in D\). La función tiene un \emph{máximo local} en
\(c\) si existe un entorno \(U\) de \(c\) tal que

\[
  f(x)\le f(c) \qquad(x\in D\cap U).
\]

El \emph{mínimo local} se define invirtiendo la desigualdad. Un máximo o
mínimo es \emph{absoluto} en \(D\) si la comparación vale para todo \(x\in D\).
El punto \(c\) es el punto extremizador; \(f(c)\) es el valor extremo.

Usaremos la siguiente convención:

\begin{itemize}
  \item un \emph{punto estacionario} es un punto interior donde \(f'\) existe
        y vale cero;
  \item un \emph{punto crítico} es un punto interior donde \(f'(c)=0\) o donde
        la derivada no existe.
\end{itemize}

Los extremos de un intervalo cerrado se consideran candidatos absolutos, pero
no se denominan críticos bajo esta convención. El teorema de Fermat de la
Sección~\ref{sec:rolle-valor-medio} afirma que todo extremo local interior en
el que \(f\) sea derivable es estacionario. Es una condición necesaria, no
suficiente: \(x^3\) tiene un punto estacionario en cero y no tiene allí un
extremo.

\subsection{Criterio de la primera derivada}
\CanonicalUnit{CM-C04-U007}{}{}
\ProofRole{proof-026}{}

\begin{teorema}[Cambio de signo de la primera derivada]
Sea \(f\) continua en un entorno de \(c\) y derivable a ambos lados de \(c\),
salvo quizá en \(c\). Entonces:

\begin{itemize}
  \item si \(f'>0\) a la izquierda de \(c\) y \(f'<0\) a la derecha, \(f\)
        tiene un máximo local en \(c\);
  \item si \(f'<0\) a la izquierda de \(c\) y \(f'>0\) a la derecha, \(f\)
        tiene un mínimo local en \(c\);
  \item si \(f'\) conserva el signo al atravesar \(c\), este criterio no
        produce un extremo local.
\end{itemize}
\end{teorema}

La justificación no es una tabla formal: el signo de \(f'\) y el teorema del
valor medio establecen la monotonía en cada intervalo lateral. Crecer hasta
\(c\) y decrecer después fuerza un máximo; decrecer y luego crecer fuerza un
mínimo.

\subsection{Criterio de la segunda derivada}
\CanonicalUnit{CM-C04-U008}{}{}
\ProofRole{proof-026}{}

\begin{teorema}[Criterio de la segunda derivada]
Supongamos que \(f'(c)=0\) y que \(f''(c)\) existe. Si \(f''(c)>0\), \(f\)
tiene un mínimo local estricto en \(c\). Si \(f''(c)<0\), tiene un máximo
local estricto en \(c\).
\end{teorema}

\begin{proof}[Derivación del criterio]
Si \(f''(c)>0\), por la definición de derivada aplicada a \(f'\), para \(x\)
suficientemente próximo a \(c\) se cumple

\[
  \frac{f'(x)-f'(c)}{x-c}>0.
\]

Como \(f'(c)=0\), resulta \(f'(x)<0\) a la izquierda y \(f'(x)>0\) a la
derecha. El criterio de la primera derivada da un mínimo. El caso
\(f''(c)<0\) es análogo.
\end{proof}

Si \(f''(c)=0\), el criterio es inconcluso. En \(c=0\), \(x^4\) tiene un
mínimo, \(-x^4\) tiene un máximo y \(x^3\) no tiene ninguno; en los tres casos
se anulan la primera y la segunda derivadas. Los criterios basados en la
primera derivada no nula de orden superior requieren hipótesis y herramientas
adicionales; no se usarán como mecanismo automático en el núcleo de este
capítulo.

\subsection{Extremos absolutos en un intervalo cerrado}
\CanonicalUnit{CM-C04-U009}{}{}
\ProofRole{proof-015}{}

Si \(f:[a,b]\to\R\) es continua, el teorema del valor extremo de la
Sección~\ref{sec:teoremas-existencia} garantiza que alcanza máximo y mínimo
absolutos. Ese resultado prueba \emph{existencia}; la derivada ayuda a
\emph{localizarlos}.

El procedimiento lógico es:

\begin{enumerate}
  \item identificar los puntos críticos del interior;
  \item incluir los extremos \(a\) y \(b\);
  \item evaluar \(f\) en todos los candidatos;
  \item comparar los valores alcanzados.
\end{enumerate}

Si la función no es continua en todo el intervalo, o el dominio no es cerrado
y acotado, la existencia debe justificarse de otra manera. Resolver \(f'=0\)
sin comprobar esos aspectos solo produce una lista parcial de candidatos.

\section{Optimización en una variable}
\label{sec:optimizacion-variable}

\subsection{Del enunciado al modelo}
\CanonicalUnit{CM-C04-U010}{}{}

Un problema de optimización aplicada no empieza derivando. Primero hay que
construir y delimitar el modelo:

\begin{enumerate}
  \item identificar las variables y sus unidades;
  \item traducir las restricciones;
  \item elegir la función objetivo;
  \item reducir el problema a una variable;
  \item determinar el dominio factible;
  \item justificar si el óptimo absoluto debe existir;
  \item localizar puntos críticos y extremos del dominio;
  \item comparar candidatos e interpretar el resultado en el contexto.
\end{enumerate}

La función objetivo expresa qué magnitud se maximiza o minimiza; una
restricción expresa qué combinaciones son admisibles. Confundirlas puede dar
una derivación correcta para un modelo equivocado.

\begin{ejemplo}[Rectángulo de perímetro fijado]
Sea \(P>0\) un perímetro dado. Si los lados de un rectángulo miden \(x\) e
\(y\), la restricción es

\[
  2x+2y=P,
\]

y el objetivo es maximizar el área \(A=xy\). Al escribir
\(y=P/2-x\), el dominio factible es \(0\le x\le P/2\) y

\[
  A(x)=x\left(\frac P2-x\right).
\]

La continuidad en el intervalo cerrado garantiza la existencia de máximo y
mínimo absolutos. Como

\[
  A'(x)=\frac P2-2x,
  \qquad
  A''(x)=-2,
\]

el único punto estacionario es \(x=P/4\), donde hay un máximo local. Los
extremos factibles dan área cero, mientras que

\[
  A\left(\frac P4\right)=\frac{P^2}{16}.
\]

Por tanto, el rectángulo de área máxima es el cuadrado de lado \(P/4\). La
respuesta final incluye la geometría y las unidades: no es solo el número
\(P/4\).
\end{ejemplo}

\section{Convexidad, concavidad e inflexión}
\label{sec:convexidad-inflexion}

\subsection{Definición mediante cuerdas}
\CanonicalUnit{CM-C04-U011}{}{}
\ProofRole{proof-027}{}

\begin{definicion}[Función convexa]
Sea \(I\) un intervalo. Una función \(f:I\to\R\) es \emph{convexa} si, para
todo \(x,y\in I\) y todo \(t\in[0,1]\),

\[
  f\bigl((1-t)x+ty\bigr)
  \le (1-t)f(x)+tf(y).
\]
\end{definicion}

El lado derecho describe la cuerda que une \((x,f(x))\) con \((y,f(y))\): la
gráfica de una función convexa queda por debajo de sus cuerdas. La función es
\emph{estrictamente convexa} si la desigualdad es estricta cuando
\(x\ne y\) y \(0<t<1\). Una función es \emph{cóncava} si \(-f\) es convexa;
en ese caso se invierten las desigualdades.

Esta es la convención canónica. No se usarán como términos formales las
expresiones históricas ``convexa hacia arriba'' o ``convexa hacia abajo''.

Una función convexa finita es continua en el interior de su intervalo, pero
la definición por sí sola no fuerza continuidad en un extremo incluido. Por
ejemplo, la función que vale cero en \([0,1)\) y vale uno en \(1\) es convexa
en \([0,1]\) y no es continua en el extremo derecho.

\subsection{Caracterizaciones diferenciales}
\CanonicalUnit{CM-C04-U012}{}{}
\ProofRole{proof-027}{}

\begin{teorema}[Criterios de convexidad]
Sea \(I\) un intervalo abierto y \(f:I\to\R\) derivable. Son equivalentes:

\begin{enumerate}
  \item \(f\) es convexa en \(I\);
  \item \(f'\) es no decreciente en \(I\);
  \item para todos \(x,y\in I\),
        \[
          f(y)\ge f(x)+f'(x)(y-x).
        \]
\end{enumerate}

Así, las rectas tangentes de una función convexa quedan por debajo de la
gráfica. Si además \(f\) es dos veces derivable en \(I\), estas condiciones
equivalen a

\[
  f''(x)\ge0 \qquad(x\in I).
\]
\end{teorema}

\begin{proof}[Derivación orientada]
Para una función convexa, las pendientes de las cuerdas aumentan al desplazar
sus extremos hacia la derecha. Al hacer tender ambos extremos hacia un punto,
las pendientes límite muestran que \(f'\) es no decreciente y producen la
desigualdad de la tangente. Recíprocamente, esa desigualdad aplicada en \(x\)
y en \(y\), con pesos adecuados, da la desigualdad de convexidad.

Si \(f''\ge0\), el criterio de monotonía aplicado a \(f'\) muestra que \(f'\)
es no decreciente. En sentido inverso, si \(f'\) es no decreciente y es
derivable, entonces \(f''\ge0\) en cada punto. Para concavidad se invierten
todos los signos.
\end{proof}

\subsection{Convexidad estricta y optimización}
\CanonicalUnit{CM-C04-U013}{}{}

Si \(f''(x)>0\) en un intervalo, entonces \(f'\) es estrictamente creciente y
\(f\) es estrictamente convexa. Es una condición suficiente, no necesaria:
\(x^4\) es estrictamente convexa en \(\R\), aunque \(f''(0)=0\).

Si \(f\) es derivable y convexa, todo punto estacionario \(c\) es un mínimo
absoluto, porque la desigualdad de la tangente da

\[
  f(y)\ge f(c)+f'(c)(y-c)=f(c)
  \qquad(y\in I).
\]

Si la función es estrictamente convexa, ese mínimo es único. Esta observación
es útil en optimización, pero no elimina la necesidad de comprobar el dominio
factible ni la existencia del candidato.

\subsection{Puntos de inflexión}
\CanonicalUnit{CM-C04-U014}{}{}

\begin{definicion}[Punto de inflexión]
Un punto interior \(c\) del dominio es un \emph{punto de inflexión} si la
función cambia de convexa a cóncava, o de cóncava a convexa, al atravesar
\(c\).
\end{definicion}

El criterio operativo consiste en estudiar la convexidad en intervalos a
ambos lados. Si \(f''\) existe y cambia de signo en \(c\), hay inflexión. La
igualdad \(f''(c)=0\) solo identifica un candidato y no es suficiente.

La función \(x^3\) cambia de cóncava en \(( -\infty,0)\) a convexa en
\((0,+\infty)\), de modo que \(0\) es un punto de inflexión. En cambio,
\(x^4\) es convexa a ambos lados de cero: aunque \(f''(0)=0\), no hay
inflexión. Tampoco es necesario que exista \(f''(c)\); lo decisivo es el
cambio de convexidad.

Los criterios de orden superior no se aplicarán sin declarar el primer orden
no nulo, la posición interior y la regularidad local requeridos. Su tratamiento
general se deja como enriquecimiento posterior a la aproximación de Taylor.

\subsection{Enriquecimiento: propiedad de Darboux de las derivadas}
\CanonicalUnit{CM-C04-U015}{}{}

Aunque una derivada no sea continua, posee la propiedad del valor intermedio.
Si \(f\) es derivable en un intervalo abierto \(I\), \(x_1<x_2\) y
\(\lambda\) está entre \(f'(x_1)\) y \(f'(x_2)\), existe
\(c\in(x_1,x_2)\) tal que \(f'(c)=\lambda\).

La estrategia considera \(g(x)=f(x)-\lambda x\). Según la orientación de los
signos de \(g'(x_1)\) y \(g'(x_2)\), la función alcanza un mínimo o un máximo
local interior; el teorema de Fermat da \(g'(c)=0\). La formulación cubre ambos
casos y no supone un orden particular entre los valores extremos de la
derivada. Este resultado no se convierte aquí en una familia de ejercicios.
